% =====================================================================
\chapter{Procesos de decisión de Markov y el LLM como política}
\label{tema:1}
% =====================================================================

Este tema introduce el lenguaje del aprendizaje por refuerzo: agente, entorno, recompensa, retorno, política y funciones de valor. La presentación sigue el capítulo~3 de Sutton y Barto~\cite{sb} y la clase~2 del curso de David Silver~\cite{silver}; la última sección traslada el formalismo a la generación de texto, como en la clase~17 de CS336~\cite{cs336}.

\section{La interfaz agente-entorno}

En cada instante $t=0,1,2,\dots$ un \emph{agente} observa el estado $S_t$ del \emph{entorno}, elige una acción $A_t$ y recibe, en el instante siguiente, una recompensa numérica $R_{t+1}$ y un nuevo estado $S_{t+1}$. La sucesión resultante
\[
S_0,\,A_0,\,R_1,\,S_1,\,A_1,\,R_2,\,S_2,\,\dots
\]
se llama \emph{trayectoria}. Cuando la interacción termina en un estado terminal tras un número finito de pasos, la trayectoria se llama \emph{episodio}.

\begin{definicion}[Proceso de decisión de Markov finito]\label{def:mdp}
Un \emph{proceso de decisión de Markov} (MDP) finito es una tupla $(\Sset,\Aset,p,\rho_0,\gamma)$ formada por un conjunto finito de estados $\Sset$, un conjunto finito de acciones $\Aset$, una \emph{dinámica}
\[
p(s',r\mid s,a)=P\bigl(S_{t+1}=s',\,R_{t+1}=r\mid S_t=s,\,A_t=a\bigr),
\]
que para cada par $(s,a)$ es una distribución sobre $\Sset\times\R$ con soporte finito, una distribución inicial $\rho_0$ sobre $\Sset$ y un \emph{factor de descuento} $\gamma\in[0,1]$. La palabra «Markov» expresa que la distribución de $(S_{t+1},R_{t+1})$ depende de la historia solo a través de $(S_t,A_t)$.
\end{definicion}

\begin{definicion}[Política]\label{def:politica}
Una \emph{política} (estacionaria y estocástica) es una aplicación que asigna a cada estado $s$ una distribución $\pi(\cdot\mid s)$ sobre $\Aset$. Una política \emph{parametrizada} $\pi_\theta$ depende de un vector de parámetros $\theta\in\R^{d}$.
\end{definicion}

Fijados un MDP y una política, la probabilidad de un prefijo de trayectoria $\tau=(s_0,a_0,r_1,s_1,\dots,s_T)$ se obtiene por la regla de la cadena (módulo~1):
\begin{equation}\label{eq:ptau}
p_\pi(\tau)=\rho_0(s_0)\prod_{t=0}^{T-1}\pi(a_t\mid s_t)\,p(s_{t+1},r_{t+1}\mid s_t,a_t).
\end{equation}
Esta factorización separa lo que controla el agente (los factores $\pi$) de lo que impone el entorno (los factores $p$), y es la base del teorema del gradiente de política del \cref{tema:2}.

\section{Recompensa y retorno}

El objetivo del agente no es maximizar la recompensa inmediata sino una cantidad acumulada~\cite[cap.~3]{sb}.

\begin{definicion}[Retorno]\label{def:retorno}
El \emph{retorno} desde el instante $t$ es
\[
G_t=R_{t+1}+\gamma R_{t+2}+\gamma^{2}R_{t+3}+\dots=\sum_{k=0}^{\infty}\gamma^{k}R_{t+k+1}.
\]
En una tarea episódica de horizonte $T$ se toma $R_k=0$ para $k>T$, de modo que la suma es finita y se admite $\gamma=1$.
\end{definicion}

\begin{proposicion}\label{prop:retorno}
\begin{enumerate}
  \item Para todo $t$, $G_t=R_{t+1}+\gamma\,G_{t+1}$.
  \item Si $\abs{R_k}\le R_{\max}$ para todo $k$ y $\gamma<1$, entonces $\abs{G_t}\le R_{\max}/(1-\gamma)$.
\end{enumerate}
\end{proposicion}

La recursión (i) permite calcular todos los retornos de un episodio en una sola pasada hacia atrás, empezando por $G_T=0$ (\cref{eje:1-retornos}).

\section{Funciones de valor y ecuación de Bellman}

\begin{definicion}[Funciones de valor y ventaja]\label{def:valor}
Sea $\pi$ una política. La \emph{función de valor de estado} y la \emph{función de valor de acción} son
\[
V^{\pi}(s)=\E_{\pi}\bigl[G_t\mid S_t=s\bigr],\qquad Q^{\pi}(s,a)=\E_{\pi}\bigl[G_t\mid S_t=s,\,A_t=a\bigr],
\]
donde $\E_\pi$ indica que las acciones futuras se eligen con $\pi$. La \emph{función ventaja} es $A^{\pi}(s,a)=Q^{\pi}(s,a)-V^{\pi}(s)$.
\end{definicion}

La ventaja mide cuánto mejor (o peor) que la media de la propia política es elegir $a$ en $s$ y continuar después con $\pi$. Es la cantidad que pondera los gradientes de todos los métodos de los temas~2, 3 y~6.

\begin{teorema}[Ecuación de Bellman para $V^{\pi}$]\label{teo:bellman}
Para todo estado $s$,
\[
V^{\pi}(s)=\sum_{a}\pi(a\mid s)\,Q^{\pi}(s,a),\qquad Q^{\pi}(s,a)=\sum_{s',r}p(s',r\mid s,a)\bigl[r+\gamma V^{\pi}(s')\bigr].
\]
En particular, $V^{\pi}(s)=\sum_{a}\pi(a\mid s)\sum_{s',r}p(s',r\mid s,a)\bigl[r+\gamma V^{\pi}(s')\bigr]$.
\end{teorema}

\begin{corolario}\label{cor:ventaja-media}
Para todo estado $s$, $\displaystyle\sum_{a}\pi(a\mid s)\,A^{\pi}(s,a)=0$.
\end{corolario}

La ecuación de Bellman es un sistema lineal de $\abs{\Sset}$ ecuaciones con $\abs{\Sset}$ incógnitas. Para $\gamma<1$ tiene solución única (\cref{eje:1-contraccion}); el capítulo~4 de~\cite{sb} la resuelve de forma iterativa.

\section{El modelo de lenguaje como política}

La generación autorregresiva encaja en la \cref{def:mdp} con una dinámica especialmente sencilla~\cite{cs336}.

\begin{definicion}[MDP de tokens]\label{def:mdp-tokens}
Sean $\V$ un vocabulario con un símbolo de fin y $x$ una instrucción. El \emph{MDP de tokens} asociado a $x$ tiene como estados las secuencias $s_t=(x,y_1,\dots,y_t)$, como acciones los tokens $a_t=y_{t+1}\in\V$, una transición determinista $s_{t+1}=(s_t,a_t)$ (concatenación), estado inicial $s_0=x$ y estados terminales los que acaban en el símbolo de fin o alcanzan una longitud máxima. La recompensa es nula salvo en la transición final, en la que vale $r(x,y)$ para la respuesta completa $y=y_{1:T}$, y se toma $\gamma=1$.
\end{definicion}

\begin{proposicion}\label{prop:mdp-tokens}
En el MDP de tokens, con una política $\pi_\theta(y_{t+1}\mid x,y_{1:t})$:
\begin{enumerate}
  \item la probabilidad de una trayectoria coincide con la de la respuesta, $p_{\pi_\theta}(\tau)=\pi_\theta(y\mid x)=\prod_{t=1}^{T}\pi_\theta(y_t\mid x,y_{<t})$;
  \item el retorno desde cualquier instante es $G_t=r(x,y)$, y $V^{\pi_\theta}(x)=\E_{y\sim\pi_\theta(\cdot\mid x)}[r(x,y)]$;
  \item para un prefijo $s_t$ y un token $a$, $Q^{\pi_\theta}(s_t,a)=V^{\pi_\theta}(s_t a)$, donde $s_t a$ denota la concatenación.
\end{enumerate}
\end{proposicion}

\begin{table}[h]
\centering\small
\begin{tabular}{@{}ll@{}}
\toprule
Elemento del MDP & En un modelo de lenguaje \\
\midrule
Estado $s_t$ & instrucción $x$ más los tokens generados $y_{1:t}$ \\
Acción $a_t$ & siguiente token $y_{t+1}$; $\abs{\Aset}=\abs{\V}$, del orden de $10^{4}$--$10^{5}$ \\
Transición & determinista y conocida: concatenar el token \\
Política $\pi_\theta(a\mid s)$ & softmax de los logits del Transformer (módulo~14) \\
Recompensa & normalmente solo al final: $r(x,y)$ \\
Episodio & una respuesta completa, hasta el símbolo de fin \\
\bottomrule
\end{tabular}
\caption{Correspondencia entre los elementos de un MDP y la generación de texto.}
\label{tab:mdp-llm}
\end{table}

\addtocounter{figure}{1}%

\begin{observacion}[Bandido contextual]\label{obs:bandido}
Como la transición es conocida y la recompensa llega solo al final, la respuesta entera $y$ puede verse como una única acción tomada en el contexto $x$ con probabilidad $\pi_\theta(y\mid x)$. Este caso particular de MDP de un solo paso se llama \emph{bandido contextual}. Los métodos de preferencias de los temas~4 y~5 trabajan a este nivel; PPO (\cref{tema:3}) trabaja token a token con ventajas $A^{\pi}(s_t,a_t)$. El número de acciones del bandido es enorme: con $\abs{\V}=5\cdot 10^{4}$ y $T=100$ hay $\abs{\V}^{T}\approx 10^{470}$ respuestas, por lo que ninguna suma sobre $y$ puede calcularse explícitamente.
\end{observacion}

\addtocounter{definicion}{1}%

